\subsection{Advantages and Limitations of SSD Technology}

\subsubsection{Performance Advantages}
\begin{enumerate}[label=\roman*.]
    \item \textbf{Low Access Time:} Data is accessed electronically without seek time or rotational latency \cite{openai2026nvmearchitecture}.
    \item \textbf{High Read/Write Speed:} Sequential and random data transfers are significantly faster than HDDs \cite{openai2026nvmearchitecture}.
    \item \textbf{Fast Boot and Application Loading:} Operating systems and applications launch much more quickly due to reduced storage latency \cite{gfg_storage_structure}.
\end{enumerate}

\subsubsection{Reliability and Durability}

\begin{enumerate}[label=\roman*.]
    \item \textbf{Shock Resistance:} With no moving components, SSDs are much more resistant to vibration and accidental drops \cite{lexar_nvme}.
    \item \textbf{Silent Operation:} SSDs generate virtually no mechanical noise \cite{lexar_nvme}.
    \item \textbf{Lower Heat Generation:} Reduced mechanical activity results in lower operating temperatures \cite{lexar_nvme}.
\end{enumerate}

\subsubsection{Limitations}
Despite their advantages, SSDs also have several limitations \cite{gfg_storage_structure}.

\begin{enumerate}[label=\roman*.]
    \item \textbf{Limited Program/Erase Cycles:} NAND flash memory cells gradually wear out after repeated write and erase operations \cite{openai2026nvmearchitecture}.
    \item \textbf{Higher Cost per Gigabyte:} SSDs generally remain more expensive than HDDs for the same storage capacity \cite{gfg_storage_structure}.
    \item \textbf{Performance Degradation During Heavy Writes:} Sustained write workloads may reduce write speed until background maintenance operations are completed \cite{openai2026nvmearchitecture}.
    \item \textbf{Complex Data Recovery:} Recovering deleted or damaged data from SSDs is considerably more difficult than from traditional HDDs due to flash memory management techniques \cite{lexar_nvme}.
\end{enumerate}